\documentclass[11pt,a4paper]{article}
\usepackage{vespa_style}

\begin{document}

\vespatitlepage
  {Difficultés techniques}
  {Verrous levés et verrous reportés}
  {Document de clôture — juillet 2026\\
   Rédigé rétrospectivement à partir du code, des archives de conversation,\\
   des fichiers de configuration du Jetson et des relevés de banc.}

\tableofcontents
\newpage

% =====================================================================
\section{Comment lire ce document}

Ce document recense les obstacles \textbf{non triviaux} rencontrés au cours du projet
\vespa{} — ceux dont la solution n'était ni évidente ni documentée, et ceux qui n'ont pas
été résolus.

Il ne s'agit pas d'un journal de bord. Un bug de compilation n'y figure pas ; un bug qui a
coûté une semaine, faussé une mesure ou contraint une décision d'architecture, oui.

\begin{note}[Convention]
\textbf{\textcolor{vespagreen}{LEVÉ}} — le verrou a été franchi. La cause racine et le
correctif sont donnés, et le correctif est dans le dépôt.

\textbf{\textcolor{vespared}{REPORTÉ}} — le verrou est ouvert. L'état exact des tentatives
est donné, ainsi que la piste jugée la plus prometteuse. \textbf{C'est là que doit
commencer toute reprise du projet.}
\end{note}

\begin{avertissement}[À lire avant de reprendre le matériel]
Plusieurs de ces verrous ne sont \emph{pas visibles dans le code}. Ils sont dans le
silicium, dans le brochage de la carte, ou dans un fichier de configuration du noyau.
Un repreneur qui lirait uniquement les sources rencontrerait à nouveau, une par une,
les mêmes impasses. C'est précisément la raison d'être de ce document.
\end{avertissement}

Une remarque transverse, et elle est structurante. Ce projet a été développé avec une
assistance forte de modèles de langage. Cette assistance a été décisive sur le volume de
code produit. Mais \textbf{sur les verrous matériels, elle a été prise en défaut de manière
répétée}, avec une grande assurance de façade. Les erreurs de diagnostic les plus coûteuses
du projet sont des hypothèses d'IA plausibles, bien argumentées, et fausses — qui n'ont été
écartées que par la mesure. La section~\ref{sec:ia} en fait le bilan explicite, car c'est
une leçon de méthode réutilisable.

% =====================================================================
\newpage
\section{Verrous levés}

% ---------------------------------------------------------------------
\subsection{Faire cohabiter deux caméras MIPI à adresse I\textsuperscript{2}C identique}

\begin{verroulevé}{Double flux OV9281 simultané}
\textbf{Symptôme.} Les deux modules Arducam OV9281 portent la \emph{même} adresse
I\textsuperscript{2}C, figée en dur dans le silicium du capteur. La documentation
courante conclut qu'un multiplexeur I\textsuperscript{2}C matériel est nécessaire pour en
exploiter deux simultanément.

\textbf{Cause racine.} L'adresse est bien unique par capteur, mais les deux caméras ne
partagent pas le même bus : chaque port CSI du Jetson dispose de son propre bus
I\textsuperscript{2}C. Le conflit d'adresse est donc \emph{apparent}, pas réel — à
condition que le noyau instancie bien un sous-périphérique par port.

\textbf{Correctif.} Noyau Arducam personnalisé + \emph{overlay} d'arbre de périphériques
dédié au mode double :
\fichier{tegra234-p3767-camera-p3768-arducam-dual.dtbo}, chargé au démarrage
(voir \fichier{jetson/extlinux.conf.reference}). Les deux capteurs apparaissent alors comme
deux n\oe uds \fichier{/dev/video*} indépendants, chacun avec son propre sous-périphérique
de contrôle. Aucun multiplexeur matériel n'est nécessaire.
\end{verroulevé}

C'est le résultat que le rapport de soutenance qualifiait d'« inédit ». La formulation est
excessive — la solution est une conséquence de l'architecture du Tegra — mais elle a bien
levé un blocage réputé rédhibitoire, et elle a économisé un composant.

% ---------------------------------------------------------------------
\subsection{Synchroniser les deux capteurs à quelques microsecondes près}

\begin{verroulevé}{Déclenchement matériel commun par PWM du Tegra}
\textbf{Symptôme.} La triangulation stéréo est fausse dès que les deux images ne sont pas
prises au même instant. Sur une cible mobile, un déphasage de quelques millisecondes
suffit à déplacer le point 3D reconstruit de plusieurs centimètres. Or Linux n'est pas un
RTOS : un déclenchement logiciel des deux caméras est soumis à l'ordonnanceur et dérive.

\textbf{Cause racine.} Toute synchronisation pilotée par le CPU est structurellement
inadéquate. Il faut sortir la génération d'horloge du domaine logiciel.

\textbf{Correctif.} Les \textbf{deux} caméras sont forcées en mode esclave
(\fichier{trigger\_mode=1}) et reçoivent le \emph{même} front de déclenchement, généré par
le \textbf{PWM matériel du Tegra} — indépendant de la charge CPU :

\begin{lstlisting}[language=bash]
/sys/class/pwm/pwmchip3/pwm0/period     = 16666666  ns  -> 60 Hz
/sys/class/pwm/pwmchip3/pwm0/duty_cycle =  1666666  ns  -> 10 %
\end{lstlisting}

La simultanéité est ainsi vraie \emph{par construction} : les deux capteurs exposent sur le
même front électrique. Implémentation : \fichier{HardwareSyncController.cpp}
(\texttt{armTrigger}, \texttt{forceSlaveMode}).

\textbf{Résultat mesuré.} Dérive inter-caméras de \textbf{0,000 ms}. Un garde-fou logiciel
subsiste dans \texttt{captureSynchronizedPair()} : toute paire dont les horodatages
s'écartent de plus de 5~ms est rejetée et resynchronisée par consommation de la file la
plus en retard.
\end{verroulevé}

% ---------------------------------------------------------------------
\subsection{Les réglages optiques ne s'appliquaient pas}

\begin{verroulevé}{Fenêtre d'application des contrôles V4L2}
\textbf{Symptôme.} Les réglages d'exposition et de gain écrits dans le JSON étaient
silencieusement ignorés par le capteur. Aucune erreur, aucun code de retour :
simplement, l'image ne changeait pas.

\textbf{Cause racine — double.}
\begin{enumerate}[nosep]
  \item \textbf{Fenêtre temporelle.} Les registres actifs du capteur ne « écoutent » pas
        tant que ses PLL ne sont pas stabilisées. Pousser les réglages avant que le flux
        ne tourne revient à écrire dans le vide.
  \item \textbf{Nom de contrôle.} Le pilote V4L2 attend strictement
        \fichier{analogue\_gain} (orthographe britannique). Toute variante
        (\fichier{gain}, \fichier{analog\_gain}) échoue \emph{en silence}. L'exposition
        fonctionnait, le gain non — ce qui rendait le diagnostic contre-intuitif.
\end{enumerate}

\textbf{Correctif.} Séquence d'initialisation stricte dans \texttt{initializeRig()} :
démarrer les flux $\rightarrow$ \textbf{attendre 500~ms} $\rightarrow$ charger le JSON
$\rightarrow$ appliquer $\rightarrow$ forcer le mode esclave $\rightarrow$ purger les files.
Le ticket de dette technique correspondant est \textbf{clos}.
\end{verroulevé}

% ---------------------------------------------------------------------
\subsection{La console série se taisait dès que le potentiomètre était branché}

\begin{verroulevé}{Conflit \broche{PA2} — potentiomètre \emph{vs} USART2\_TX}
\textbf{Symptôme.} Sur le \fichier{motion\_node}, brancher le potentiomètre de réglage fin
de l'axe \textsc{pan} rendait la console série muette ou illisible. Aucun lien apparent
entre les deux.

\textbf{Cause racine.} Le shield X-NUCLEO-IHM16M1 câble son entrée « SPEED »
(potentiomètre) sur \broche{PA2}. Mais sur la carte NUCLEO-G431RB, \broche{PA2} est
\textbf{déjà câblée sur \texttt{USART2\_TX}}, c'est-à-dire sur le port COM virtuel du
ST-LINK. Y injecter une tension analogique revient à court-circuiter la sortie série.

\textbf{Correctif.} Potentiomètre \textsc{pan} déplacé sur \broche{PB14}
(\fichier{pinout.h}). Le fichier historique \fichier{STM32\_WiringManifest.txt} indique
encore \broche{PA2} : \textbf{il est faux sur ce point}, et
\fichier{docs/annexes/cablage.md} le corrige.
\end{verroulevé}

% ---------------------------------------------------------------------
\subsection{Le SPI des encodeurs pouvait mourir sans raison}

\begin{verroulevé}{\broche{PA4} : NSS matériel \emph{vs} retour de courant}
\textbf{Cause racine.} \broche{PA4} est simultanément la broche \texttt{SENSEW} du moteur 1
(retour de courant phase W) et le \textbf{NSS matériel du SPI1}. Or SPI1 porte les deux
encodeurs AS5048A, dont dépend tout l'asservissement. Un NSS matériel laissé flottant peut
faire basculer le périphérique SPI en mode esclave et interrompre la lecture des encodeurs.

\textbf{Correctif — un arbitrage, pas une solution.} \broche{PA4} est maintenue haute et
réservée au SPI. \textbf{Le retour de courant de la phase W du moteur 1 est sacrifié.}
C'est acceptable ici : la FOC de SimpleFOC reconstruit la phase manquante
($i_W = -i_U - i_V$), et surtout l'asservissement actuel est en position sur encodeur, sans
boucle de courant. \textbf{Un repreneur qui voudrait activer la mesure de courant devra
rouvrir cet arbitrage.}
\end{verroulevé}

% ---------------------------------------------------------------------
\subsection{Deux programmes se disputaient le bus I\textsuperscript{2}C}

\begin{verroulevé}{Verrou d'exclusion mutuelle matériel}
\textbf{Symptôme.} Lancer un outil Python de calibration pendant que le n\oe ud C++ tournait
corrompait les réglages du capteur, voire figeait le flux.

\textbf{Cause racine.} Les deux processus écrivent sur les mêmes sous-périphériques
I\textsuperscript{2}C. Rien, au niveau du système, ne les en empêche.

\textbf{Correctif.} Verrou exclusif \texttt{flock} sur \fichier{/tmp/vespa\_hardware.lock},
pris au démarrage de \texttt{HardwareSyncController} (\texttt{acquireHardwareLock()}).
Le second processus \textbf{refuse de démarrer} avec un message explicite.

\textbf{Ce n'est pas un bug :} lancer le n\oe ud de vision et un outil Python simultanément
échoue \emph{volontairement}.
\end{verroulevé}

% ---------------------------------------------------------------------
\subsection{\texttt{jetson-io.py} détruisait la configuration des caméras}

\begin{verroulevé}{Abandon de l'outil NVIDIA au profit d'un arbre de périphériques maîtrisé}
\textbf{Symptôme.} Impossible d'avoir \emph{en même temps} le pilote caméra et le PWM
matériel. Configurer l'un via \fichier{jetson-io.py} cassait l'autre.

\textbf{Cause racine.} \fichier{jetson-io.py} régénère et réapplique les \emph{overlays}
d'arbre de périphériques de façon globale : il écrase les configurations de multiplexage
existantes au lieu de les compléter.

\textbf{Correctif.} L'outil est \textbf{proscrit}. La configuration matérielle est gérée
par des \emph{overlays} explicites, versionnés, chargés depuis
\fichier{/boot/extlinux/extlinux.conf}. Les fichiers sont conservés dans
\fichier{jetson/} (voir \fichier{docs/annexes/jetson\_setup.md}).

\begin{avertissement}[Ne pas lancer \texttt{jetson-io.py}]
Sur une machine reconstruite selon ce dépôt, exécuter \fichier{jetson-io.py} et sauvegarder
\textbf{cassera la synchronisation stéréo}. L'\emph{overlay} PWM
(\fichier{jetson-io-hdr40-user-custom.dtbo}, broches 32 et 33) doit être préservé tel quel.
\end{avertissement}
\end{verroulevé}

% ---------------------------------------------------------------------
\subsection{Valider le protocole CAN sans bus CAN}

\begin{verroulevé}{Bus virtuel \texttt{vcan0}}
Le bus physique n'ayant jamais fonctionné (voir §\ref{sec:can}), la totalité de la couche
protocolaire a été validée sur une interface CAN \textbf{virtuelle} :

\begin{lstlisting}[language=bash]
sudo modprobe vcan
sudo ip link add dev vcan0 type vcan && sudo ip link set up vcan0
candump vcan0
\end{lstlisting}

Les trames \texttt{TARGET\_VEC} (\texttt{0x020}, 7 octets) et \texttt{FIRE\_CMD}
(\texttt{0x010}) ont été observées, correctement sérialisées, à la cadence attendue.
\textbf{La sérialisation, les identifiants, les priorités et la logique applicative sont
donc démontrés corrects} ; seule la couche physique est en défaut. C'est une distinction
importante pour le repreneur : \textbf{il n'y a pas de bug à chercher dans le protocole.}
\end{verroulevé}

% =====================================================================
\newpage
\section{Verrous reportés}

% ---------------------------------------------------------------------
\subsection{Le bus CAN physique : les broches n'ont jamais été réclamées}
\label{sec:can}

\begin{verroureporte}{CAN0 du Jetson — \texttt{pinmux} non réclamé}
\textbf{C'est le verrou bloquant du projet.} Les deux n\oe uds fonctionnent
individuellement. Ils n'ont jamais communiqué \emph{sur du cuivre}.
\end{verroureporte}

\paragraph{Symptôme.} L'interface \fichier{can0} n'apparaît jamais. La vérification du
multiplexeur ne renvoie \textbf{rien} :

\begin{lstlisting}[language=bash]
$ sudo bash -c "cat /sys/kernel/debug/pinctrl/*/pinmux-pins | grep -i paa"
$          # <- aucune sortie : les broches PAA ne sont meme pas listees
\end{lstlisting}

\paragraph{Le piège de brochage.} Sur le Tegra 234 :
\begin{center}
\begin{tabular}{@{}lll@{}}
\toprule
\textbf{Broches SoC} & \textbf{Contrôleur} & \textbf{En-tête 40 broches} \\
\midrule
\broche{PAA0} / \broche{PAA1} & \textbf{CAN0} & broche 31 (TX) / broche 29 (RX) \\
\broche{PAA2} / \broche{PAA3} & CAN1 & — \\
\bottomrule
\end{tabular}
\end{center}

\noindent
Assigner la fonction \texttt{can0} aux broches \broche{PAA2}/\broche{PAA3} revient à
demander au silicium de router \emph{CAN0 sur les broches physiques de CAN1}. Le pilote
\texttt{pinctrl} détecte l'incohérence et \textbf{s'arrête}, ce qui explique que les broches
ne soient plus listées du tout.

\paragraph{État final laissé sur la machine.} L'\emph{overlay} chargé au démarrage est
\fichier{/boot/can0-pinmux.dtbo}, dont l'en-tête interne annonce
« \texttt{CAN0 Hardware Override V5} » — cinquième itération. Il déclare bien les n\oe uds
corrects (\texttt{can0\_dout\_paa0}, \texttt{can0\_din\_paa1}, fonction \texttt{can0}) et
tente de neutraliser les n\oe uds \texttt{paa2}/\texttt{paa3}. \textbf{Il n'a pas suffi.}

\begin{avertissement}[Une hypothèse répandue — et fausse]
La dernière piste explorée avant l'arrêt du projet attribuait la panne à un \emph{overlay}
\fichier{jetson-io-hdr40-user-custom.dtbo} « empoisonné », qui aurait contenu une mauvaise
affectation CAN sur \broche{PAA2}/\broche{PAA3}.

\textbf{Cette hypothèse est démentie par le fichier lui-même.} Son inspection (juillet 2026)
montre qu'il ne configure \emph{que} deux broches :
\texttt{hdr40-pin32}~$\rightarrow$~\texttt{soc\_gpio19\_pg6} et
\texttt{hdr40-pin33}~$\rightarrow$~\texttt{soc\_gpio21\_ph0} — les broches du PWM de
synchronisation. \textbf{Il ne contient aucune configuration CAN.}

Ne pas perdre de temps à « réparer » ce fichier : il n'est pas en cause, et le modifier
casserait la synchronisation stéréo. La cause est ailleurs.
\end{avertissement}

\paragraph{Piste recommandée pour la reprise.} Par ordre de coût croissant :
\begin{enumerate}
  \item \textbf{Vérifier le domaine de \emph{pinmux}.} Les broches \broche{PAA} du Tegra
        appartiennent au contrôleur \textbf{AON} (\emph{always-on}), distinct du
        \emph{pinmux} principal. Un \emph{overlay} qui cible le mauvais contrôleur n'a
        aucun effet — et c'est cohérent avec l'absence totale des broches dans
        \texttt{debugfs}. \textbf{C'est la piste la plus probable.} Vérifier que
        l'\emph{overlay} cible bien \texttt{pinmux\_aon} et non \texttt{pinmux}.
  \item \textbf{Confirmer que le module noyau \texttt{mttcan} se charge}
        (\texttt{modprobe mttcan}, puis \texttt{dmesg | grep -i can}). Le contrôleur CAN du
        Tegra en dépend.
  \item \textbf{Contourner le SoC.} Solution de repli assumée et rapide : un contrôleur CAN
        \textbf{externe sur SPI} (MCP2515 ou équivalent). On abandonne le CAN natif du
        Tegra, mais on obtient un bus fonctionnel sans livrer bataille au \emph{pinmux}.
        \textbf{Si la reprise est sous contrainte de temps, commencer par là.}
\end{enumerate}

\paragraph{Ce que ce verrou n'est pas.} Il n'est ni logiciel, ni protocolaire. Le câblage
(WCMCU-230 / SN65HVD230, 3,3~V, broches 29/31, masse commune, terminaison 120~$\Omega$) est
correct et documenté. Le code est validé sur \texttt{vcan0}. \textbf{Seul le routage des
broches est en cause.}

% ---------------------------------------------------------------------
\newpage
\subsection{Reconnaissance biologique : VespAI est en couleur, nos capteurs non}

\begin{verroureporte}{Incompatibilité VespAI / capteur monochrome}
\textbf{Le système ne sait pas reconnaître un frelon.} Il suit un marqueur ArUco.
\end{verroureporte}

\paragraph{L'enchaînement des contraintes.} Ce verrou n'est pas une erreur : c'est une
cascade logique dont chaque maillon est justifié, et dont le résultat est une impasse.

\begin{enumerate}
  \item Le flou de mouvement détruit la détection $\rightarrow$ il faut un
        \textbf{obturateur global}.
  \item Aucun capteur \textbf{trichrome} à obturateur global n'existe à un prix accessible
        $\rightarrow$ le capteur sera \textbf{monochrome} (OV9281).
  \item Le modèle VespAI~\cite{osheawheller2024vespai} est entraîné sur des images
        \textbf{RGB} $\rightarrow$ \textbf{incompatible}.
\end{enumerate}

\paragraph{Ce qui a été tenté.} Une « chirurgie réseau » sur les poids pré-entraînés, sans
réentraînement. L'idée est mathématiquement exacte : si l'on duplique le canal gris sur les
trois canaux RGB, la première convolution calcule
\[
  \text{sortie} = I_R W_R + I_G W_G + I_B W_B
  = I_{\text{gris}}\,(W_R + W_G + W_B),
\]
donc \textbf{sommer les poids de la première couche sur la dimension canal} produit un
modèle à 1 canal \emph{strictement équivalent}, sans perte, sans jeu de données, sans
entraînement.

\begin{lstlisting}[language=Python]
new_weights = torch.sum(old_weights, dim=1, keepdim=True)  # (out,3,k,k) -> (out,1,k,k)
first_conv_layer.weight      = nn.Parameter(new_weights)
first_conv_layer.in_channels = 1
\end{lstlisting}

\paragraph{Pourquoi cela a échoué quand même.} La transformation est exacte, et le modèle
obtenu tourne vite (compilé en TensorRT FP16, entrée native
$1{,}280 \times 800$ — les deux dimensions étant multiples de 32, aucun redimensionnement
n'est requis, donc pas de recours au VIC). \textbf{Mais l'équivalence mathématique ne crée
pas d'information.} Le modèle reste entraîné sur des statistiques de couleur : il détecte des
frelons \emph{imprimés sur papier blanc}, et échoue en conditions réelles, avec des faux
positifs massifs. Surtout, la distinction \emph{V.~velutina} / \emph{V.~crabro} — l'exigence
éthique centrale du projet — repose largement sur la \textbf{coloration}
\cite{sipos2025morphologie} : elle est \textbf{structurellement impossible} en monochrome à
partir d'un modèle couleur.

\paragraph{Décision.} À trois semaines de la soutenance, le réentraînement complet
(estimé 1,5 à 2,5 semaines, sans garantie de résultat, le jeu de données n'étant pas public)
a été jugé trop risqué. \textbf{Repli assumé sur un marqueur ArUco}
\cite{garridojurado2014aruco} : haut contraste, natif en niveaux de gris, aucun
entraînement, pose 3D exacte. Le marqueur \textbf{tient lieu de cible} et permet de valider
toute la chaîne aval (triangulation, prédiction, asservissement) — qui est, elle,
indépendante de la nature de la cible.

\begin{note}[Pour la reprise]
La voie n'est pas la « chirurgie réseau ». C'est la \textbf{constitution d'un jeu de données
monochrome} (images issues de l'OV9281 lui-même, en conditions réelles) et le
réentraînement d'un détecteur dessus. Tant que la discrimination inter-espèces ne sera pas
rétablie, \textbf{le système ne doit pas tirer} : il ne sait pas distinguer une espèce
invasive d'une espèce protégée.
\end{note}

% ---------------------------------------------------------------------
\subsection{Le n\oe ud de coordination, le watchdog, et le tir}

\begin{verroureporte}{Watchdog non implémenté — laser jamais mis sous tension}
Le \fichier{coordination\_node} (ESP32) \textbf{n'existe pas}. Il est spécifié — ses trames
CAN sont définies (\texttt{FIRE\_CMD}, \texttt{POS\_COORD}, \texttt{HB\_COORD}) — mais
aucune ligne de code ne l'implémente.
\end{verroureporte}

Or ce n\oe ud porte trois fonctions, dont une est une condition de sécurité :

\begin{center}
\begin{tabularx}{\textwidth}{@{}lX@{}}
\toprule
\textbf{Fonction} & \textbf{État} \\
\midrule
Collimation dynamique & Loi de commande établie, mécanisme spécifié, jamais assemblé \\
Séquence de tir & Non implémentée \\
\textbf{Watchdog} & \textbf{Non implémenté} — \emph{condition sine qua non du tir} \\
\bottomrule
\end{tabularx}
\end{center}

\paragraph{Conséquence, et elle est délibérée.} Le watchdog est le mécanisme qui surveille
les \emph{heartbeats} de tous les n\oe uds et coupe l'alimentation générale en cas de
défaut ou de silence. Sans lui, \textbf{rien ne garantit l'extinction du laser si un n\oe ud
se fige}.

\begin{avertissement}[La diode laser n'a jamais été alimentée]
Ce n'est pas un retard : c'est un \textbf{refus}. Mettre sous tension une diode de
\textbf{Classe 4} (voir \fichier{docs/annexes/securite\_laser.md}) sans \emph{interlock}
fonctionnel aurait été une faute. Aucun tir n'a été effectué, et \textbf{aucun ne doit
l'être} tant que le watchdog n'est pas opérationnel et testé.
\end{avertissement}

\paragraph{Note de conception.} Le watchdog est prévu \emph{dans le n\oe ud de tir} —
c'est-à-dire au plus près de l'organe dangereux, sur un c\oe ur CPU dédié. C'est le bon
choix : un superviseur qui dépendrait du réseau pour couper le courant serait vulnérable à
la panne même qu'il doit traiter. La coupure est \textbf{matérielle} (alimentation
générale), l'inhibition logicielle du tir n'intervenant que pendant le temps de décharge des
condensateurs.

% ---------------------------------------------------------------------
\subsection{La profondeur de champ létale : un résultat inattendu}

\begin{verroureporte}{Collimation dynamique — jamais assemblée ni étalonnée}
\end{verroureporte}

Le mécanisme de collimation dynamique vise à \textbf{dégrader volontairement} le faisceau
partout ailleurs qu'au point d'impact, afin de réduire la dangerosité intrinsèque du
dispositif. Deux lentilles : \textbf{L1} (fournie avec le module, rapprochée de la source
pour la faire \emph{diverger}) et \textbf{L2} (longue focale, mobile, qui refocalise à une
distance commandée).

La loi de commande a été établie analytiquement par les relations de conjugaison de
Descartes, en éliminant la focale inconnue $f_1$ au profit de grandeurs mesurables
(demi-divergence $\theta$, diamètre $D_{L2}$ du faisceau sur la lentille mobile) :
\[
  x_{\text{cible}} \;=\; x_{L2} \;+\;
  \frac{f_2 \cdot \dfrac{D_{L2}}{2\tan\theta}}
       {\dfrac{D_{L2}}{2\tan\theta} - f_2}
\]
Cette relation, quadratique, est approximée par un polynôme à $10^{-4}$~mm près — plus
rapide à évaluer sur microcontrôleur.

\paragraph{Le résultat qu'il faut retenir.} Le tracé de densité de puissance
(\fichier{simulation/optique/Laser\_Optics\_v1.m}) donne deux enseignements
\textbf{opposés}, et c'est le second qui compte :

\begin{itemize}
  \item \textcolor{vespagreen}{\textbf{Bonne nouvelle.}} Le seuil d'embrasement de matières
        sèches n'est franchi que sur $\approx 50$~mm de part et d'autre de la cible. Un tir
        manqué venant frapper le sol présente donc un risque d'incendie
        \textbf{limité}. L'objectif du mécanisme est atteint.
  \item \textcolor{vespared}{\textbf{Mauvaise nouvelle, et elle domine.}} Le seuil de
        dangerosité \textbf{oculaire} reste dépassé sur près de \textbf{200~mètres}. La
        collimation dynamique \textbf{ne rend pas le dispositif sûr} : elle réduit le risque
        d'incendie, pas le risque oculaire.
\end{itemize}

\begin{avertissement}
Il serait dangereux de conclure de ce mécanisme que le système est « moins dangereux ».
Il ne l'est pas pour l'\oe il. La collimation dynamique est une mesure de mitigation
\textbf{partielle}, en aucun cas un substitut aux protections (lunettes, \emph{interlock},
confinement).
\end{avertissement}

% ---------------------------------------------------------------------
\subsection{Verrous mineurs et dette technique}

\begin{center}
\begin{tabularx}{\textwidth}{@{}p{4.2cm}Xl@{}}
\toprule
\textbf{Sujet} & \textbf{Description} & \textbf{État} \\
\midrule
Appels \texttt{std::system} &
Le pilotage du PWM passe par des appels shell (\texttt{echo > /sys/class/pwm/...}).
Coûteux et fragile. À remplacer par des écritures \texttt{sysfs} directes en C++. &
Ouvert \\
\addlinespace
Chemins absolus &
Les \texttt{\#define} de scripts PWM sont des chemins absolus : déplacer le dépôt casse la
compilation. À rendre relatifs ou configurables. &
Ouvert \\
\addlinespace
Échelle d'exposition &
L'OV9281 n'utilise pas le même calcul de temps-ligne en mode maître et en mode
déclenchement externe. L'outil \fichier{app\_EG.py} calibre donc dans un régime différent
de celui du \emph{pipeline} C++ : les valeurs ne correspondent pas 1:1. &
Ouvert \\
\addlinespace
Timing des contrôles V4L2 &
Attendre que le flux tourne avant de pousser les réglages. &
\textcolor{vespagreen}{\textbf{Clos}} \\
\addlinespace
Suivi multi-cibles &
Un seul filtre de Kalman. Une nuée de frelons n'est pas gérée : ni association de données,
ni priorisation. &
Ouvert \\
\bottomrule
\end{tabularx}
\end{center}

% =====================================================================
\newpage
\section{Ce que l'assistance par IA a coûté, et ce qu'elle a rapporté}
\label{sec:ia}

Cette section est inhabituelle dans un document technique. Elle est ici parce que le projet
a été, de bout en bout, co-développé avec des modèles de langage, et que \textbf{le bilan
n'est pas uniforme}. Un repreneur qui adopterait la même méthode doit savoir où elle porte
et où elle nuit.

\paragraph{Ce qui a très bien fonctionné.} La production de code. Le volume de logiciel écrit
— chaîne V4L2 zéro-copie, sérialisation CAN, FOC, triangulation, filtre de Kalman, outils
web de calibration — est hors de portée de trois étudiants sur la durée impartie sans cette
assistance. Le rapport de soutenance le dit, et c'est exact.

\paragraph{Ce qui a coûté cher.} \textbf{Le diagnostic matériel.} Les modèles produisent des
hypothèses plausibles, argumentées, confiantes — et fausses. Elles sont d'autant plus
coûteuses qu'elles sont \emph{crédibles} : elles envoient chercher au mauvais endroit,
parfois pendant des jours.

\begin{center}
\begin{tabularx}{\textwidth}{@{}p{5.2cm}X@{}}
\toprule
\textbf{Affirmation de l'IA} & \textbf{Réalité établie par la mesure} \\
\midrule
« Deux OV9281 exigent un multiplexeur I\textsuperscript{2}C matériel » &
Faux. Chaque port CSI a son propre bus. Un \emph{overlay} adapté suffit. \\
\addlinespace
« \fichier{jetson-io.py} : sélectionnez \texttt{can0}, cela routera les broches » &
Faux, et destructeur : l'outil écrase la configuration caméra/PWM. \\
\addlinespace
« L'\emph{overlay} \fichier{jetson-io-hdr40-user-custom.dtbo} est empoisonné : il assigne
CAN0 à \broche{PAA2}/\broche{PAA3} » &
\textbf{Faux.} Inspection du binaire : il ne contient \emph{que} les deux broches PWM
(32 et 33). Aucune configuration CAN. \textbf{Le projet s'est arrêté sur cette fausse
piste.} \\
\addlinespace
« La chirurgie réseau donnera un détecteur équivalent » &
Mathématiquement exact, \textbf{opérationnellement faux} : l'équivalence formelle ne
compense pas l'absence d'information couleur. \\
\bottomrule
\end{tabularx}
\end{center}

\paragraph{La règle qui s'en dégage.} Sur ce projet, \textbf{toutes} les avancées matérielles
décisives sont venues d'une \emph{mesure} qui contredisait une hypothèse — pas d'un
raisonnement. Le conflit \broche{PA2}/série a été trouvé en débranchant. Le
\emph{pinmux} CAN a été caractérisé en interrogeant \texttt{debugfs}, pas en raisonnant sur
des registres. La fausse piste finale aurait été écartée en \emph{ouvrant le fichier
incriminé} — ce qui prend deux minutes, et n'a pas été fait.

\begin{note}[Recommandation au repreneur]
Utilisez ces outils pour écrire du code : le gain est réel et considérable. Mais devant un
symptôme matériel, \textbf{n'acceptez aucune cause racine qui n'ait été confirmée par une
mesure}. Le coût d'une vérification est presque toujours inférieur au coût d'une fausse
piste bien argumentée.
\end{note}

% =====================================================================
\newpage
\section{Synthèse : par où reprendre}

\begin{center}
\begin{tabularx}{\textwidth}{@{}clX@{}}
\toprule
\textbf{\#} & \textbf{Chantier} & \textbf{Pourquoi dans cet ordre} \\
\midrule
1 & \textbf{Débloquer le CAN physique} &
Rien d'intégré n'est démontrable tant que les n\oe uds ne se parlent pas. Piste : domaine
\emph{pinmux} \textbf{AON}. Repli rapide : contrôleur CAN externe sur SPI. \\
\addlinespace
2 & \textbf{Implémenter le watchdog} &
Condition \emph{sine qua non} de toute mise sous tension du laser de Classe 4.
Aucun essai de tir avant cela. \\
\addlinespace
3 & \textbf{Jeu de données monochrome + réentraînement} &
Rend au système sa fonction biologique et la discrimination inter-espèces. Sans cela, le
système ne doit pas tirer. \\
\addlinespace
4 & \textbf{Assembler et étalonner la collimation} &
Loi de commande établie ; mécanisme jamais monté. Attention : ne réduit
\emph{pas} le risque oculaire. \\
\addlinespace
5 & \textbf{Gestion multi-cibles} &
Un filtre de Kalman par cible + association de données. \\
\bottomrule
\end{tabularx}
\end{center}

\vspace{1cm}
\begin{note}[Le mot de la fin]
Le système, tel qu'il est laissé, \textbf{suit} une cible en trois dimensions, à 60~Hz, avec
une prédiction qui compense sa propre latence, et il oriente deux axes asservis en boucle
fermée pour la garder en visée. Ce qui manque n'est pas le c\oe ur du problème : c'est
\textbf{un bus qui ne s'allume pas}, \textbf{un jeu de données qui n'existe pas}, et
\textbf{un garde-fou qu'il fallait écrire avant d'oser tirer}.

Les fondations sont saines. Les trois chantiers sont indépendants. Bonne reprise.
\end{note}

% =====================================================================
\newpage
\bibliographystyle{plain}
\bibliography{vespa}

\end{document}
